\documentclass{NSF}

% ----------------------------------------------
% Draw a line halfway down the page
% See https://tex.stackexchange.com/questions/281424/draw-a-line-halfway-down-a-page
% 
% \usepackage{eso-pic}
% \AddToShipoutPictureBG{%
%     \AtTextCenter{\hspace{-0.5\textwidth}\rule{\textwidth}{0.5pt}}}
% 
% ----------------------------------------------

\begin{document}

\clearpage
\setcounter{page}{1}

% ===============================
\section*{Appendix 3: Equipment}
% \marginpar{\color{red} 0.5 pages}
% ===============================

% ----------------------------
\subsection{Cornell University (Marohn laboratory except where noted)}\hfill
% ----------------------------

\begin{itemize}

\item Two custom-built, vacuum electric force microscopes.  Unique capabilities include CW and 0.5~ns pulsed laser illumination, variable-temperature operation, broad-band local dielectric spectroscopy (quantitative charge density, conductivity, and dielectric constant); phase-kick electric force microscopy (transient conductivity, nanosecond resolution); and fast turn-around (5 to 10 samples/day).

\item Lead-halide perovskite film preparation: double MBraun UNIlab glove box with Laurel WS-650MZ spin coater equipped with custom automated gas quench system, Trovato glove box with evaporator, UV-Ozone substrate cleaner, SCS 6800 Spin Coater.

%These microscopes employ:
%
%    \begin{itemize}
%    \item a fast-scanning stage enabling the analysis of three samples in one pump-down;
%    \item two turbo pumps for fast turn around, enabling 5 to 10 samples per day;
%    \item custom-built fiber-optic Fabry-Perot interferometers with 50~fm/$\sqrt{\mathrm{Hz}}$ displacement sensitivity;
%    \item Minus-k Technologies vibration isolation stages; and 
%    \item RHK PLLPro Universal AFM Controllers.
%    \end{itemize}

\item Zeiss Gemini/Sigma for scanning electron microscopy and EDX (Cornell CCMR facility).

\item Bruker D8 GADDS for X-ray diffraction (Cornell CCMR facility).

\end{itemize}

% ----------------------------
\subsection{National Laboratory of the Rockies (NLR)}\hfill
% ----------------------------

\begin{itemize}

\item  Femtosecond photoluminescence, transient absorbance, and transient terahertz spectroscopies.
\item Structure characterization using X-ray diffraction.  Surface and interface characterization via Auger electron spectroscopy, X-ray photoelectron spectroscopy, and secondary ion mass spectrometry.
\item AIPS, an Automated In-situ Photodosing Spectrometer allowing continuous absorption spectra to be taken at approximately 1 spectra per minute per sample contained in a sealed environmental chamber for complete control of atmosphere, light intensity up to 1 sun, and temperature.
\item Steady state microwave conductivity with illumination for dark and photoconductivity, outfitted with an automated X stage for up to 10 samples including automated data collection and processing.
\item Transient microwave conductivity. 
QYB LuQY Pro PLQE measurement system, designed specifically for measuring PLQE and QFLS in perovskite materials and devices.

\end{itemize}

\end{document}